\hypertarget{appendix-deca299ae4b8}{}
\section*{ULP distance policies}
\section{ULP distance policies}\label{ulp-distance-policies}

\textbf{Curriculum depth:} Introduced · \textbf{First seen:} C004 · \textbf{Latest development:} C004

\subsection{Definition and mechanism}\label{definition-and-mechanism}

Distance counts representable steps under a declared ordering/equality policy. The finite-only laboratory collapses both zeros to one key and rejects NaNs/infinities. Across negative minimum subnormal to positive minimum subnormal distance is two.

\subsection{Prerequisites}\label{prerequisites}

\begin{itemize}
\tightlist
\item
  \hyperref[appendix-14bab4cab9d4]{ULP and representable spacing} --- Adjacent representable points define rounding candidates, errors and step distances.
\item
  \hyperref[appendix-3a048f032c1f]{Signed zero} --- Numerical equality merges two encodings and must be reflected in the chosen ordering or byte policy.
\item
  \hyperref[appendix-d6607cc4c235]{NaNs and payloads} --- Unordered comparisons and payload distinctions require explicit metric, comparison and serialization policy.
\end{itemize}

\subsection{Used by}\label{used-by}

Connect later dependent concepts here as their lessons arrive.

\subsection{Recall and next depth}\label{recall-and-next-depth}

Reconstruct the concrete example in the definition and identify its validity assumptions.

\textbf{Pending:} A raw complement key with two zero slots overcounts cross-zero intervals unless corrected.

\subsection{Source and navigation}\label{source-and-navigation}

\href{Continuous\%20Notes/Computing\%20Foundations\%20-\%20From\%20Bits\%20to\%20C++\%20and\%20Systems\%20APIs\#A\%20stretching\%20ruler\%20for\%20real-valued\%20quantities}{M004 teaching chapter} · \hyperref[appendix-fd090d982a15]{Per-part definitions} · \href{}{Dependency Map (02-Foundations-II, appendix-4ad47bdd269a)}
